\chapter{Joins and Query Optimization}
\label{ch:joins}

A SPARQL query of $k$ triple patterns is a $k$-way join, and its running time
depends far more on \emph{how} it is joined than on anything else. This chapter
covers the three join algorithms, how an optimizer estimates the size of
intermediate results, and how it searches for a good join order.

\section{Three join algorithms}

Let the outer input have $|R|$ rows and the inner input $|S|$ rows.

\paragraph{Nested-loop join.} For each outer row, scan the whole inner input and
keep the pairs whose shared slots agree. Cost: $|R| \times |S|$ comparisons, no
extra memory. Correct for any join, and the fallback when nothing better applies.

\paragraph{Index nested-loop join.} For each outer row, substitute its bindings
into the inner pattern and \emph{probe an index} for just the matching triples.
If the outer row binds the subject and SPO exists, the inner side becomes a
prefix range, often one or a handful of triples. Cost: $|R|$ index probes. This
is the common, fast case in SPARQL: one pattern finds the entities, and each
further pattern looks up an attribute of each entity.

\paragraph{Hash join.} \emph{Build} a hash table on the shared slots of one input
(ideally the smaller), then \emph{probe} it with each row of the other. Cost:
linear in $|R| + |S|$, plus memory for the build side. It wins when both inputs
are large and no index makes the probe selective, and it loses when the build
side does not fit in memory, or when an index probe was already a point lookup.

\begin{keyidea}
Index nested loop is fast when the outer side is small and the inner side is
indexed on the join key; hash join is fast when both sides are large; nested
loop is always correct. An optimizer's job is to know which case it is in.
\end{keyidea}

\section{Cardinality estimation}

To compare plans, the optimizer needs the \term{cardinality} of each
intermediate result: how many rows each operator produces. It estimates them
from \term{statistics} gathered at load time: the total triple count $|T|$,
the number of distinct subjects, predicates, and objects, and per-predicate
counts.

For a pattern with predicate $p$ and object $o$ bound, the classic estimate
assumes the positions are \term{independent}:
\[
  \widehat{\,|\{(s,p,o)\}|\,} \;=\; \frac{|p| \cdot |o|}{|T|},
\]
where $|p|$ and $|o|$ are the counts of triples with that predicate and that
object. Independence is convenient and frequently wrong. In RDF, predicates and
objects are strongly correlated (\code{rdf:type} goes with class objects; a
\code{name} predicate never has a class as its object), so independence can be
off by orders of magnitude.

The \term{q-error} measures an estimate $\hat c$ against the true count $c$:
\[
  q(\hat c, c) = \max\!\left(\frac{\hat c}{c},\ \frac{c}{\hat c}\right) \ \ge 1 .
\]
It treats over- and underestimates symmetrically, and errors multiply as they
propagate up a plan: a join of two inputs each off by 10 can be off by 100.

\section{Searching for a join order}

A left-deep plan over $k$ patterns has $k!$ orders, and each join can use
several algorithms. Selinger's System R optimizer \cite{selinger} made the
search tractable with dynamic programming:

\begin{lstlisting}[language=pseudo]
for each single pattern P:   best[{P}] = cheapest access path for P
for size = 2 to k:
  for each subset S of patterns with |S| = size:
    for each pattern P in S:
      candidate = join(best[S - {P}], P) with the cheapest algorithm
      if cost(candidate) < cost(best[S]): best[S] = candidate
return best[all patterns]
\end{lstlisting}

It keeps one best plan per subset rather than per order, which reduces $k!$ to
about $k \cdot 2^k$. A simpler \term{greedy} strategy (start with the most
selective pattern, then repeatedly add the cheapest connected pattern) is a
good first step and often good enough for the star-shaped joins common in SPARQL.

\section{Explain analyze}

The single most useful optimizer tool is to \emph{run} the plan and print, next
to each operator's estimated rows, the rows it actually produced. The node where
the two diverge most is where a statistic is missing, usually a correlation that
the independence assumption averaged away.

\begin{ontodb}[: the optimizer]
OntoDB's optimizer takes the planner's tree and returns another tree that must
return the same rows. The shell setting \code{\textbackslash set join
auto|nlj|inlj|hash} either lets it choose or forces one algorithm, and forcing
happens in the optimizer, not the executor, so the plan that
\code{\textbackslash explain} prints is always the plan that runs. Statistics
are collected at load: the triple count, distinct counts per position,
per-predicate counts, and counts for partially bound patterns. Estimates are
finite, never negative, and binding one more position never increases them.
Join ordering starts greedy and left-deep.
\end{ontodb}

\begin{decision}{Exact counts for single patterns}
\dpart{Choice} Statistics are collected in one scan of the store: counts per
subject, predicate, and object, per pair of positions, and the set of triples.
Every single triple pattern of constants therefore has an \emph{exact}
cardinality; estimation, and its independence assumptions, begin only at joins.
\dpart{Alternatives} Distinct counts and per-predicate counts only, combined
with the independence formula above; histograms or samples; or RDF-3X's
characteristic sets, which summarize which predicates occur together on one
subject.
\dpart{Why} The independence formula is wrong by orders of magnitude on RDF,
where predicates and objects are correlated. For a store of this size, exact
pair counts are affordable and remove that error from every scan estimate.
\dpart{Cost} About seven hash entries per triple, collected at load, and
stale after updates until the next collection, like a table that has not been
re-\code{ANALYZE}d.
\end{decision}

\begin{decision}{Join estimates use distinct values}
\dpart{Choice} A join's estimate is System R's
$|L|\,|R| / \max(V(L,x), V(R,x))$ for each shared variable $x$, where $V$ is
the number of distinct values $x$ takes. For a single pattern, $V$ comes from
statistics: a predicate's distinct subjects or objects, or the pattern's own
count when the other positions are bound.
\dpart{Alternatives} Assume each side's join values are all distinct, which
gives $\min(|L|, |R|)$ for one shared variable; or a fixed selectivity per join.
\dpart{Why} SPARQL joins come in two shapes that the simple rule cannot both
get right. On a store of 100 students taking 3 of 10 courses, joining students
to their names (a star on the subject) is 100 rows, and joining enrolments to
the course class is 300. Distinct values give both exactly; $\min(|L|,|R|)$
says 10 for the second, 30 times too low.
\dpart{Cost} Two more counts per predicate, and an assumption that remains: the
join values of the smaller side are taken to appear on the larger side. A
variable's $V$ in an intermediate result is approximated by the smaller of its
$V$ in a pattern and the result's size.
\end{decision}

\begin{decision}{The inner side is re-opened with the outer row}
\dpart{Choice} Every executor has \code{InitWith(row)}, which by default is
\code{Init}. A triple scan overrides it: each variable that the outer row
already binds becomes a constant in the pattern it hands the store. The nested
loop re-opens its inner side this way for every outer row.
\dpart{Alternatives} Open the inner scan once with its own pattern and compare
every inner row with every outer row; or materialize the inner side once and
loop over the copy.
\dpart{Why} With the outer row filled in, the inner pattern usually binds the
subject, so on an indexed store each probe is a short prefix range instead of
a scan of every triple with that predicate. The cost per outer row drops from
the size of the inner pattern to the size of its match.
\dpart{Cost} One store scan opened, and one tree descent, per outer row, with
no caching between rows that bind the same values.
\end{decision}

\begin{decision}{Index nested loop probes the store directly}
\dpart{Choice} The index nested-loop join requires its inner side to be one
triple pattern. It has no inner executor: for each outer row it builds the
probe pattern from the outer row and asks the store, which picks the index
whose order starts with the bound positions.
\dpart{Alternatives} Treat any nested loop over a rebound scan as an index
nested loop, since with rebinding it already probes.
\dpart{Why} With rebinding, the nested loop already probes an index, so the
two operators read the same triples. The separate operator skips the inner
executor and the generic row merge, and it names the strategy in the plan,
where \code{\textbackslash explain} and the optimizer can see it.
\dpart{Cost} Only a scan can be the inner side; a join whose inner side is
itself a join stays a nested loop.
\end{decision}

\begin{decision}{Hash join builds on the inner side}
\dpart{Choice} The hash join reads its inner (right) input once into a table
keyed by the ids of the variables both sides produce, then streams the outer
input through it. A row whose key slot is unbound is not hashed but compared
with every row of the other side.
\dpart{Alternatives} Build on whichever side is smaller, or partition both
sides to disk when the table does not fit (a grace hash join).
\dpart{Why} Choosing the smaller side needs cardinality estimates, which
arrive with the optimizer (plans 5.1 and 5.2); until then the inner side of a
left-deep plan is a single pattern, usually the smaller input.
\dpart{Cost} The whole inner side lives in memory; a large inner pattern can
make the build expensive or exhaust memory. There is no spilling.
\end{decision}

\begin{decision}{Forcing a join method rewrites the plan}
\dpart{Choice} \code{\textbackslash set join hash} (or \code{nlj}, \code{inlj})
is applied by the optimizer (\code{Optimizer::ApplyJoinMethod}), which swaps
each nested-loop node for the chosen one before the plan is explained or run.
The executor factory only reads node types.
\dpart{Alternatives} Keep one join node and let the executor read the setting
at run time.
\dpart{Why} \code{\textbackslash explain} then always prints the plan that
runs. A test can check that a forced hash plan contains a hash join without
running it.
\dpart{Cost} The optimizer must be invoked even when the user forces
everything.
\end{decision}

\begin{decision}{Greedy, left-deep join ordering first (planned)}
\dpart{Choice} Start from the most selective pattern and repeatedly add the
cheapest pattern connected to what is already joined, producing a left-deep
tree.
\dpart{Alternatives} Selinger's dynamic programming over subsets, or a search
that also considers bushy trees.
\dpart{Why} Most SPARQL basic graph patterns are stars around one subject, for
which greedy is close to optimal, and a greedy order is easy to explain in
\code{\textbackslash explain analyze}.
\dpart{Cost} Chains and snowflakes, where a bushy plan or a non-greedy first
step wins, can be ordered badly.
\end{decision}

\begin{status}
Today the planner always joins in textual order with nested loop, and the
nested loop rebinds its inner scan for every outer row (plan 4.1). The index
nested-loop join (plan 4.2) and hash join (plan 4.3) run, chosen with
\code{\textbackslash set join}, and the differential test checks every query
file under all three against the memory store. Statistics (plan 5.1) and cardinality estimation (plan 5.2) run. Join
ordering, the automatic choice, and \code{\textbackslash explain analyze} are
plans 5.3 to 5.5.
\end{status}

\begin{checkyourself}
\item Outer side 1,000 rows, inner pattern 1,000,000 triples, and the join binds the inner subject with SPO available. Compare nested loop and index nested loop.
\item When does a hash join beat index nested loop?
\item With exact pair counts, what is OntoDB's estimate for \code{(?, p, o)}, and where can its estimates still go wrong?
\item A store has $|T| = 10^6$, $|p| = 10^4$, $|o| = 10^3$. What does the independence formula estimate for \code{(?, p, o)}? The true count is 1,000. What is the q-error?
\item Why do estimation errors multiply up a plan?
\item A hash join builds on \code{?c a ub:Course} (2,000 rows) and probes with \code{?s ub:takesCourse ?c} (50,000 rows). What does it keep in memory, and how many hash lookups does it make? What changes if the sides are swapped?
\item \code{?s a ub:Student . ?s ub:name ?n}: the first pattern matches 1,000 students and the second 1,000,000 triples. How many triples does the nested loop read with rebinding on an indexed store, and how many without it?
\item Why does forcing a join method in the optimizer, and not in the executor, matter for \code{\textbackslash explain}?
\end{checkyourself}

\section*{Further reading}
Selinger et al.\ \cite{selinger}; Graefe \cite{graefe}; Leis et al.\
\cite{leis2015} on how good cardinality estimates really are; Neumann and
Weikum \cite{rdf3x} on RDF selectivity.
